\section{Protocol, source, and comparator details}
\label{app:protocols}

\subsection{Exact saved-plan provenance}

The canonical K2000, Wishart, and Chook configurations are read from
the following frozen plan files in the companion \texttt{qefem\_prof}
repository:
\begingroup\footnotesize
\begin{itemize}[leftmargin=2em]
 \item \path{benchmarks/k2000/results/portable_accuracy_recheck/final_k2000/plan.json};
 \item \path{benchmarks/wishart/results/portable_accuracy_recheck/final_wishart/plan.json};
 \item \path{benchmarks/chunk/results/portable_accuracy_recheck/final_chook_a/plan.json}
       and \path{final_chook_b/plan.json}.
\end{itemize}
\endgroup
The paper workspace retains the parsed configurations and endpoint rows
under \path{Paper/results/}. It records the Python, NumPy, and PyTorch
versions, the numerical device, dtype, input hashes, and the declared
source revision in \path{results/experiment_environment.json} and the
case-level configuration and manifest files. The frozen K2000
\QeFEM{} job was selected from \path{scale_k2000/momentum_fast}.
That selection predates the later G-set manual-discrete search. No
G-set spectral or discrete-neighbour setting was silently transferred
to the plotted K2000 results.

Each K2000 seed is obtained from the declared base seed 50,000 by the
repository's deterministic derivation function. The resulting five
seed values are listed in Table~\ref{tab:k2000trials}; the LQA rows
share those labels. Wishart and Chook likewise retain their exact seed
arrays in each instance's \path{configuration.json}. A seed controls
initial pseudo-random fields, not the quality of the optimum itself.
Different algorithms map the same integer to different state shapes,
so a shared label is useful for auditing but not a proof of paired
random trajectories.

\subsection{The G-set search variants}

The smooth G-set stage used a product-state objective with exact
automatic differentiation. The later remaining-gap search used both
an exact manual evaluation of Eq.~\eqref{eq:fieldgrad} and the
surrogate direction of Eq.~\eqref{eq:discretenbr}. Its spectral schedule
computes $\kappa$ from the smallest eigenvalue of the normalized graph
coupling matrix, constructs a decreasing $q/\kappa$ curve, and chooses a
decreasing ratio $\rho=\Gamma/T$. Equation~\eqref{eq:spectralinvert}
then generates the temperature and driver arrays. Most transferred
discrete-update proposals used SGD with momentum 0.95 and weight decay
0.01; exact proposals retained their previous optimizer. This is a
graph-specific search, not one fixed set of hyperparameters across
all G-set graphs.

For the discrete mode, $\lambda_t$ can increase linearly between
declared start and end fractions or switch abruptly when those
fractions are equal. The implementation changes only the neighbor
magnetizations inside $K\widetilde{\bm m}_z$ and leaves the local
Bloch Jacobian and entropy contribution continuous. The final
configuration is always decoded and rescored against the original
graph. The full controller, per-graph plan, proposal registry, and
rescore audit are saved under
\path{benchmarks/Gset/results/remaining_discrete/}; its completed
analysis is the source of the 54-instance table in
Appendix~\ref{app:fullresults}.

\subsection{How the FEM and dSB context was handled}

The accompanying FEM repository contains a K2000 notebook with saved
cell output for separate \FEM{} and \dSB{} sweeps. The current notebook
source sets 1,000 replicas for each method and nominal update budgets
from 100 through 100,000. Its saved output reports the best cuts in
Table~\ref{tab:archived-context}. These values are included so the
reader can see that the current \QeFEM{} gap is a real competitive
limitation, not because the notebook is a controlled three-way trial.
The notebook itself states that its saved outputs predate a later
device-selection code update; rerunning is needed to align saved
values with the current MPS/CPU implementation. The published \FEM{}
paper reports its own K2000 protocol and figure
\citep{shen2025}; its figure should not be numerically identified with
this local notebook's later output.

\begin{table}[ht]
\centering\small
\caption{Selected best cuts printed by the archived FEM K2000 notebook,
with 1,000 FEM or dSB trajectories at each independent update budget.
These are contextual outputs, not matched comparisons with the separate
8,192-replica, five-trial QeFEM CPU evaluation. The stored reference is
33,337.}
\label{tab:archived-context}
\begin{tabular*}{\textwidth}{@{\extracolsep{\fill}}rrr}
\toprule
updates & archived FEM best & archived dSB best\\
\midrule
100 & 33,133 & 32,993\\
200 & 33,214 & 33,155\\
400 & 33,275 & 33,248\\
600 & 33,296 & 33,285\\
800 & 33,296 & 33,272\\
1,000 & 33,306 & 33,304\\
2,000 & 33,332 & 33,333\\
4,000 & 33,337 & 33,337\\
\bottomrule
\end{tabular*}
\end{table}

The FEM notebook uses RMSprop-like momentum settings
$\eta=0.03$, $\alpha=0.56$, momentum $0.63$, weight decay $0.013$,
$T_{\max}=1.16$, $T_{\min}=6\times10^{-5}$, and a graph-dependent
normalization in its classical FEM implementation. Its dSB implementation
uses step size $\Delta t=1.25$ and a sign-dependent coupling in a
position--momentum update. These are genuine algorithmic and protocol
differences. Figure~\ref{fig:k2000-population} consequently compares
only the frozen QeFEM population study with its explicitly declared
saved LQA baseline; it does not mix a FEM/dSB sweep into one curve.

\subsection{Figure data and rebuild command}

All four figures in this manuscript are generated by
\path{detailed/make_figures.py}. The script performs read-only loads of
saved endpoint, trajectory, and G-set analysis files; it executes no
optimizer. It writes vector PDF and PNG copies to
\path{detailed/figures/}, exact plot-ready aggregates to
\path{detailed/plot_data/}, and result-table rows to
\path{detailed/tables/}. The script asserts the expected K2000 five-seed
population counts and G-set 54-instance status total before plotting.
The figure captions identify whether they show a five-seed endpoint
summary, one coherent trajectory, fixed planted-instance means, or an
adaptive attainment inventory. Those distinctions are part of the
figure data, not small-print exceptions.
